\begin{table}[htbp]
\centering
\footnotesize
\caption{医院报告完整度对截距贡献的等效性检验，以及各变化量的配对自助法 90\% 置信区间。}
\label{tab:s3}
\begin{tabular}{llrlrlr}
\toprule
设定 & $10^3\Delta a$ [90\% CI] & $10^3\delta_{0.10}$ & 等效 & $10^3\delta_{0.05}$ & 等效 & 等效季数 \\
\midrule
$R_A$，规则窗口（$S=4$） & 0.29 [$-$0.22, 0.85] & 2.55 & 是 & 0.64 & 否 & 4/4 \\
$G$，规则窗口（$S=4$） & 0.37 [$-$2.83, 2.81] & 10.13 & 是 & 2.53 & 否 & 3/4 \\
$G$，日历窗口（$S=3$） & $-$0.38 [$-$3.55, 2.24] & 10.14 & 是 & 2.54 & 否 & 2/3 \\
\bottomrule
\end{tabular}
\par\smallskip
\begin{minipage}{0.95\linewidth}\scriptsize $\Delta a$ 为统计量的截距减去以报告医院占比校正计数后重算的统计量截距，区间来自配对双向自助法（$B=2000$）。等效界 $\delta_{\tau_0}=\kappa^2\tau_0^2(1+v_R/\bar{R}^2)$ 由报告波动阈值与各统计量的样本均值 $\bar{R}$ 确定（$G$ 取 $\kappa=1$，$v_R=0.04$）。90\% 区间落在 $(-\delta,\delta)$ 之内时判定为等效（TOST，$\alpha=0.05$）。等效季数为 $\tau_0=0.10$ 时判定等效的流行季数。\end{minipage}
\end{table}
